% Part: applied-modal-logics
% Chapter: epistemic-logic
% Section: relational-models

\documentclass[../../../include/open-logic-section]{subfiles}

\begin{document}

\olfileid{aml}{el}{rel}

\olsection{مدل‌های رابطه‌ای}

مفهوم پایهٔ معناشناختی برای منطق‌های معرفتی همان مفهوم متناظر در
منطق‌های موجهات معمولی است. مدل‌های رابطه‌ای همچنان از مجموعه‌ای از
جهان‌ها و تخصیصی تشکیل می‌شوند که تعیین می‌کند کدام !!{propositional
variable} در کدام جهان «صادق» به شمار می‌آید. اگر تنها با یک عامل سروکار
داشته باشیم، مانند معمول یک رابطهٔ دسترسی‌پذیری داریم. اما اگر منطق
معرفتی چندعاملی باشد، رابطهٔ دسترسی‌پذیریِ واحد ما به مجموعه‌ای از روابط
دسترسی‌پذیری بدل می‌شود: یک رابطه برای هر~$a$ در مجموعهٔ نمادهای
عامل~$G$.

یک \emph{مدل رابطه‌ای} از مجموعه‌ای از جهان‌ها تشکیل می‌شود که روابط
دسترسی‌پذیری دوتایی---یکی برای هر عامل---آن‌ها را به هم مرتبط می‌کنند،
به‌همراه تخصیصی که تعیین می‌کند کدام !!{propositional variable} در کدام
جهان صادق است.

\begin{defn}
  یک \emph{مدل} برای زبان معرفتی چندعاملی سه‌تاییِ
  $\mModel{M} = \tuple{W, \{R_a\}_{a\in G}, V}$ است، که در آن
  \begin{enumerate}
  \item $W$ مجموعه‌ای ناتهی از «جهان‌ها» است؛
  \item برای هر~$a \in G$، ${R}_a$ یک رابطهٔ دسترسی‌پذیری دوتایی
  روی~$W$ است؛ و
  \item $V$ تابعی است که به هر !!{propositional
    variable}~$p$ مجموعه‌ای از جهان‌های ممکن، یعنی~$V(p)$، نسبت می‌دهد.
  \end{enumerate}
  هرگاه~$R_a ww'$ برقرار باشد، می‌گوییم~$w'$ \emph{برای عامل~$a$
    از~$w$ دسترس‌پذیر} است. هرگاه~$w \in V(p)$، می‌گوییم~$p$
  \emph{در~$w$ صادق} است.
\end{defn}

سازوکار درست همان سازوکار منطق موجهات عادی است، با این تفاوت که روابط
دسترسی‌پذیری بیشتری افزوده شده‌اند. برای یک عامل معین، معمولاً رابطهٔ
دسترسی‌پذیری او را نمایانگر چیزی دربارهٔ حالت‌های اطلاعاتی‌اش تفسیر
می‌کنیم. برای نمونه، غالباً~$R_a ww'$ را چنین می‌خوانیم که~$w'$ با
اطلاعات عامل~$a$ در~$w$ سازگار است. به بیان دیگر، عامل~$a$ در~$w$
نمی‌تواند میان جهان~$w$ و جهان~$w'$ تمایز بگذارد.

\end{document}
